\clearpage
\fancyhead[L]{\small BODE. ASTRONOMISCHES JAHRBUCH.}
\hypertarget{native-630-lower}{}
\begin{center}\small Göttingische gelehrte Anzeigen. Stück 5. Seite 41..48. 1821 Januar 8.\\[3pt]\rule{.72\linewidth}{.3pt}\end{center}

\emph{Astronomisches Jahrbuch für das Jahr \textup{1823}, nebst einer Sammlung der neuesten in die astronomischen Wissenschaften einschlagenden Abhandlungen, Beobachtungen und Nachrichten. Berechnet und herausgegeben von} Dr. J. E. \textsc{Bode}, Königl. Astronom u. s. w. Berlin 1820. Bei dem Verfasser, und in Comm. bei F. Dümmler. 252 Seiten in Octav, nebst einer Kupfertafel.

Das astronomische Jahrbuch ist gegenwärtig den Freunden der Astronomie jedesmal doppelt willkommen, da es, nachdem die von den Herren von \textsc{Lindenau} und \textsc{Bohnenberger} herausgegebene astronomische Zeitschrift aufgehört hat, in Deutschland den einzigen Vereinigungspunkt für diese Wissenschaft darbietet. Die dem vorliegenden Jahrgange beigefügten Aufsätze sind folgende. Versuch über die physische Beschaffenheit der Cometen, und besonders ihres Schweifes von Hrn. Prof. \textsc{Fischer}. Man sieht hier mit Vergnügen dasjenige zusammengestellt, was die zuverlässigsten Beobachtungen und geläuterte Ansichten über diesen Gegenstand gelehrt haben. Gegen einige einzelne Behauptungen möchte sich \SourceLeaf{631} jedoch noch manches erinnern lassen. So sind bekanntlich die Meinungen der Astronomen über die Frage, ob die Cometen mit eignem Licht oder mit reflectirtem Sonnenlicht leuchten, noch getheilt, und ein um die Theorie der Cometen vielfach und hochverdienter Astronom hat in einem früheren Aufsatze so ziemlich alles, was sich nach dem gegenwärtigen Bestand unsrer Kenntnisse darüber sagen lässt, erschöpft. Der Verf. erklärt sich für die erstere Meinung, weil das Licht der kleinsten Fixsterne selbst durch die dichtern Theile der Cometenschweife dringt, und nach einigen Erfahrungen Fixsterne selbst durch den Kern gesehen sein sollen. Allein es scheint, dass er mit Unrecht hieraus schliesst, die Cometenmasse könne das Sonnenlicht gar nicht merklich reflectiren, und das Licht, damit wir die Cometen sehen, müsse nothwendig ganz deren eignes sein. Jene Erfahrung beweist nur, dass die Cometenmasse bei weitem mehr Licht durchlässt als reflectirt; aber wenn man in Erwägung zieht, wie ausserordentlich viel blässer als Planetenscheiben selbst die glänzendsten Cometen erscheinen, so hat man, auch wenn nur ein äusserst geringer Theil des Sonnenlichts zurückgeworfen wird, noch immer nicht nöthig, eignes Licht der Cometen zu Hülfe zu nehmen. Indem der Verf. annimmt, dass die Cometenmassen Stoffe enthalten, die negativ gegen die Sonne gravitiren, scheint ihm entgangen zu sein, dass die Bewegungen solcher Cometen dann nicht genau den \textsc{Kepler}'schen Gesetzen folgen würden. Inzwischen, wenn wir gleich die Möglichkeit nicht abläugnen, dass in der That das eine \textsc{Kepler}'sche Gesetz bei den Cometenbewegungen einige Modification leiden könne, und dass eigentlich dieser Umstand bei keinem Cometen methodisch untersucht sei, so lässt sich doch aus allen bisher geführten Rechnungen schliessen, dass eine solche Modification nur äusserst klein sein dürfte. --- Beobachtungen des Cometen vom Jahre 1819 auf der Sternwarte Bogenhausen bei München von Hrn. Steuerrath \textsc{Soldner}. --- Geographische Ortsbestimmungen in Ostfriesland, durch Hrn. Prof. \textsc{Oltmanns}. --- Ueber die Länge von Pisa aus astronomischen Beobachtungen von Hrn. Prof. \textsc{Wurm}. Das Resultat dieser schätzbaren Untersuchung ist, dass der \emph{Unterschied} zwischen der von \textsc{Inghirami} aus geodätischen Messungen und der aus den astronomischen Beobachtungen gefolgerten Länge wegen der schlechten Beschaffenheit der letztern eigentlich gar nichts beweisen kann. --- Von demselben Astronomen Beiträge zu geographischen Längenbestimmungen aus Beobachtungen der Sonnenfinsternisse vom 18. November 1816, und 4. Mai 1818. --- Astronomische Beobachtungen zu Wilna von Hrn. Prof. \textsc{Sniadecky}; zu Palermo von Hrn. \textsc{Cacciatore} (die Cometen von 1819 betreffend); zu Prag von den Herren Prof. \textsc{David} und \textsc{Bittner}. --- Noch etwas über den grossen Cometen und seinen Vorübergang vor der Sonne von Hrn. Doctor \textsc{Olbers}. Die schon im astronomischen Jahrbuche für 1822 angeführte Beobachtung des Hrn. General \textsc{von Lindener}, der am 26. Juni 1819 in der. Stunde, wo der Comet bestimmt vor der Sonnenscheibe stehen musste, die Sonne ganz ohne Flecken gesehen hatte, schien das, was ohnehin höchst wahrscheinlich war, vollkommen zu beweisen, dass nemlich die lockere Cometenmasse viel zu wenig Sonnenlicht auffangen könnte, um die mindeste Trübung zu verursachen. Dennoch zeigt dieser Vorfall auf eine merkwürdige Art, wie leicht in solchen Fällen, wo keine besonders geschärfte Aufmerksamkeit statt findet, Irrthum möglich ist. Mehrere andere Beobachter haben nemlich doch an demselben Tage wirkliche Sonnenflecken gesehen, der Hr. Prof. \textsc{Schumacher}, damals in Altona, und der Hr. Prof. \textsc{Brandes} in Breslau: die Beobachtungen des Hrn. \textsc{von Lindener} verliert daher freilich ihre Beweiskraft. Denn wenn die wirklichen Sonnenflecken ihm entgangen sind, so konnte er noch weniger den Cometen bemerken, wenn auch derselbe stärkern Instrumenten oder einer mehr gesteigerten Aufmerksamkeit erkennbar gewesen wäre. Und wirklich könnte man beinahe durch zwei andere Beobachtungen das letztere anzunehmen bewogen werden, wenn nicht dabei einige Nebenumstände wären, die die Sache wieder ganz ungewiss machten. Hr. Dr. \textsc{Gruithuisen} in München sah nemlich an eben diesem Tage und in der Stunde, wo der Comet vor der Sonne stand, ausser zwei andern gewöhnlichen Sonnenflecken noch einen sehr kleinen unbegrenzten, den er an \SourceLeaf{632} frühern Tagen nicht bemerkt hatte. Allein Hr. Dr. \textsc{Gruithuisen} spricht davon als von einem \emph{schwarzen} Punkte, und so konnte sich dieser Comet doch nicht zeigen, daher wir eher geneigt sein möchten, diese Erscheinung für einen gewöhnlichen erst kurz vorher entstandenen Sonnenfleck zu halten, der in den schwächern Fernröhren der Herren \textsc{Schumacher} und \textsc{Brandes} unsichtbar blieb. Ausserdem hat auch Hr. Prof. \textsc{Wildt} in Hannover um dieselbe Morgenstunde einen sehr verwaschenen Flecken in der Sonne bemerkt, ist aber rücksichtlich des Datum selbst ungewiss, da er die Nachricht erst viel später blos aus dem Gedächtniss mitgetheilt hat. Die Hauptfrage bleibt unter diesen Umständen noch immer unentschieden, und wird es vielleicht, da Conjuncturen der Art so äusserst selten sind, noch lange bleiben. --- Die Ephemeride für den Nordstern, für alle obern Culminationen des Jahrs 1821 wird allen practischen Astronomen, die die kleine Druckschrift der Herren \textsc{Struve} und \textsc{Walbeck}, woraus sie entlehnt ist, nicht selbst besitzen, sehr willkommen sein, nur hätte dieselbe wohl eine Seite mehr verdient, damit nicht die lezte Decimale hätte wegbleiben müssen. --- Astronomische Beobachtungen in Wien von Hrn. Prof. \textsc{Bürg}; in Prag von Hrn. Prof. \textsc{Hallaschka} und in Berlin von dem Hrn. Herausgeber. --- Ueber die beobachtete Existenz einer Photosphäre der Venus im Jahr 1820 von Hrn. Geh. Rath \textsc{Pastorf}. Hr. P. sah im April d.\,J. die Venus (und späterhin auch den Jupiter) mit einem kreisförmigen sehr scharf begrenzten Lichtschimmer umgeben, welchen er als eine dem Planeten selbst angehörige Lichtsphäre betrachtet. Diese Erscheinung verdient genauer untersucht zu werden, da bei Gegenständen dieser Art so leicht optische Täuschung einfliessen kann. Hr. \textsc{Pastorf} hat den Durchmesser dieses Lichtschimmers, der ein ähnliches Ansehn hat, wie die Nachtseite des Mondes im Erdlichte, im April 16 Minuten gross gefunden. Bei der von Hrn. \textsc{Pastorf} am. 27. April beobachteten Erscheinung, wo ein kleiner teleskopischer Stern oben am \emph{östlichen} Rande der Lichtsphäre der Venus eintrat, \emph{einige Minuten} unsichtbar blieb, und dann an der \emph{Westseite} wieder erschien, ist der Umstand etwas bedenklich, dass an diesem Tage das Fortrücken der Venus in jeder Minute nur 3 Raumsecunden betrug. Ueber die Hypothese des Hrn. \textsc{Pastorf}, dass die Venus einen Trabanten haben könne, der nicht über die Lichtsphäre hinauskomme, und der, wie Hr. \textsc{Pastorf} glaubt, wenn er nur einen Durchmesser von 3 Secunden habe, uns dann im reflectirten Sonnenlicht immer unsichtbar bleiben müsse, ausser dass er zuweilen wie ein dunkler Flecken vor der Venus erscheine, wollen wir dem Urtheile unsrer Leser nicht vorgreifen. --- Beobachtete gerade Aufsteigungen des Saturn und der Vesta im Jahr 1819, der Pallas und des Mars im Jahr 1820 am Mittagsfernrohr der Göttinger Sternwarte von Hrn. Hofr. \textsc{Gauss}. --- Beschreibung des auf der Königsberger Sternwarte aufgestellten \textsc{Reichenbach}'schen Meridiankreises, von Hrn. Prof. \textsc{Bessel}. --- Beobachtungen des Cometen von 1819, nebst Sternbedeckungen von Hrn. Prof. \textsc{Struve}. --- Bestimmung der Schiefe der Ekliptik mit einem \textsc{Reichenbach}'schen Meridiankreise auf der Sternwarte Bogenhausen bei München von Hrn. \textsc{Soldner}. Die hiezu angewandten Beobachtungen geben zugleich die Bestimmung des Sonnenhalbmessers, und zwar $16'\,0''93$ für die mittlere Entfernung. Wir bemerken bei dieser Gelegenheit, dass dieser Halbmesser aus 271 am \textsc{Reichenbach}'schen Mittagsfernrohr der hiesigen Sternwarte vom Januar bis Juli 1820 beobachteten Fadenantritten sich fast genau eben so gross, nemlich zu $16'\,1''01$ ergeben hat. --- Astronomische Beobachtungen zu Kremsmünster von Hrn. \textsc{Derflinger}. --- Beobachtungen der Juno, Pallas, Ceres, des Mars und Uranus, wie auch der Schiefe der Ekliptik von Hrn. Prof. \textsc{Nicolai} in Mannheim. Letztere mit dem dreifussigen \textsc{Reichenbach}'schen Repetitionskreise bestimmt ergibt sich aus dem Wintersolstitium um $6''$ kleiner, als aus dem Sommersolstitium, wodurch; so wie durch die mit demselben Instrumente bestimmten Sterndeclinationen aufs neue bestätigt wird, dass in Rücksicht auf die Einwirkung der Schwerkraft auf die Theile der astronomischen Instrumente ein jedes wie ein Individuum betrachtet werden müsse. --- Ueber die Genauigkeit der Beobachtungen am Mittagsfernrohre von Hrn. Dr. \textsc{Walbeck}. Es wird hier zum erstenmal die richtige Behandlungsart dieses \SourceLeaf{633} wichtigen Gegenstandes gelehrt und auf die trefflichen Dorpater Beobachtungen angewandt. S.\,186 Z.\,3 von unten ist durch einen Druckfehler tang $\delta$ statt sec $\delta$ gesetzt. --- Astronomische Beobachtungen zu Hamburg von Hrn. \textsc{Rümker}. --- Ueber die geographische Lage von Dresden von Hrn. Dr. \textsc{Raschig}. --- Astronomische Bemerkungen von Hrn. Prediger \textsc{Luthmer} in Hannover. --- Ueber das wahre Datum der nächtlichen Schlacht am Halys von Hrn. Prof. \textsc{Oltmanns}. In dieser trefflichen Abhandlung, von der wir ungern alle Citate und Anmerkungen der Raumersparniss wegen weggelassen sehen, wird mit grosser Evidenz gezeigt, dass die vielbesprochene, während der berühmten Schlacht eingetretene Sonnenfinsterniss keine andere gewesen sein könne, als die vom 30. September 609 vor Chr. --- Ueber die Bahn des Ponsschen (\textsc{Enke}'schen) Cometen, nebst Berechnung seines Laufs bei seiner nächsten Wiederkehr im Jahr 1822, von Hrn. Prof. \textsc{Enke}. Die wiederholte und vollständige geführte Berechnung der Störungen, welche dieser Comet von 1786--1819 erlitten hat, gibt die befriedigendste Darstellung der Beobachtungen in den vier bisherigen Erscheinungen, nur zeigt sich eine merkwürdige durch die Rechnung noch nicht zu erklärende Beschleunigung der Bewegung, indem aus den drei Umläufen von 1786--1795 eine um einen halben Tag grössere, und aus den vier Umläufen von 1805--1819 eine um einen halben Tag kleinere Umlaufszeit hervorgeht, als aus den dreien von 1795--1805. Ueber die Ursache dieses Phänomens werden sich, wenn die nächste Wiedererscheinung des Cometen ein ähnliches Resultat geben sollte, wahrscheinliche Vermuthungen angeben lassen. Für die leichtere Wiederauffindung im Jahre 1822, wo der Comet am 24. oder 25. Mai durch seine Sonnennähe gehen wird, hat Hr. Prof. \textsc{Enke} durch eine bequeme Ephemeride bestens gesorgt. Auf der südlichen Halbkugel wird die Beobachtung im Juni und Juli keine Schwierigkeiten haben; allein in Europa wird er um diese Zeit wegen seiner südlichen Lage gar nicht zu sehen sein, und früher wird seine gar zu grosse Lichtschwäche die Erkennung wenn nicht unmöglich, doch höchst schwierig machen. Doch damit nichts unversucht bleibe, wünschen wir, dass der um diesen Cometen so hoch verdiente Astronom die am 25. Februar anfangende Ephemeride noch ein paar Monat weiter rückwärts fortsetzen möge, da früher wenn gleich bei noch grösserer Lichtschwäche doch wegen des hohen Standes bei dunkler Nacht vielleicht noch etwas mehr Hoffnung statt zu finden scheint, als da, wo nach geendigter Dämmerung der Comet dem Horizont schon so nahe steht. --- Auch für den kleinern Cometen des Jahrs 1819 hat Hr. Prof. \textsc{Enke} eine elliptische Bahn mit einer Umlaufszeit von nur $5\tfrac{1}{2}$ Jahren gefunden, die sowohl die Marseiller als die erst später bekannt gewordenen Mailänder Beobachtungen vortrefflich darstellt. --- Von derselben Opposition der Vesta 1819 und Ephemeride für die nächste Erscheinung dieses Planeten. --- Ephemeriden für die Juno und Pallas 1821 von Hrn. Prof. \textsc{Nicolai}, und von Hrn. \textsc{von Staudt} in Göttingen. --- Astronomische Beobachtungen im Jahr 1820 von Hrn. Hofrath \textsc{Gauss}. --- Ueber die Bestimmung der geographischen Breite vermittelst des Polarsterns von Hrn. Prof. \textsc{Dirksen} in Berlin (Analyse einer bequemen zu diesem Zweck von Hrn. Prof. \textsc{Schumacher} gegebenen Tafel). --- Noch Beobachtungen von Sternbedeckungen, Jupiterstrabanten-Verfinsterungen und der grossen Sonnenfinsterniss vom 7. Sept. 1820 durch Hrn. Prof. \textsc{Rümker} in Hamburg, wie auch von letzterer durch Hrn. Prof. \textsc{Nicolai} und Hrn. \textsc{von Heiligenstein} in Mannheim. --- Von den am Schlusse dieses Bandes befindlichen kurzen Nachrichten zeichnen wir hier noch die von der in Abo erbaueten neuen Sternwarte und von der in den Russischen Ostsee-Provinzen vorzunehmenden Gradmessung aus. --- Die Kupfertafel stellt ausser den Sternbedeckungen und den beiden Mondfinsternissen des Jahrs 1823 noch den berechneten Vorübergang des grossen Cometen von 1819 vor der Sonne, die beobachtete geocentrische Bewegung des \textsc{Enke}'schen Cometen 1818 und 1819, und die Venus mit einem von Hrn. \textsc{Pastorff} beobachteten Flecken dar.
